\section{FOLLOW-UP AND DISCOVERY}\label{sec:design}

Auditors already have the vocabulary for the distinction the results point
at. Follow-up is returning to a previously identified problem. Discovery is
learning where problems exist outside what has already been identified. Both
are part of the work, and they are not the same job.

For follow-up, a record of prior findings is close to ideal routing
information: it identifies the known conditions that deserve revisiting, and
the regulation is right to require the visit. It is not sufficient
verification that a condition remains; the follow-up procedures generate that evidence, and the
Single Audit results are consistent with follow-up working as designed.
Discovery concerns conditions not already represented in the finding history.
Prior findings can carry genuine information about it: some entities are
durably more susceptible across the board, and one underlying weakness can
produce several kinds of condition. What the record cannot say is how much of
that apparent signal is susceptibility and how much is the observation
process, because the places without findings are precisely where absence and
non-detection cannot be separated, and because a prior finding can widen the
search. A record built by looking where the record said to look contains
information about the places that were looked at; how much it contains about
everywhere else is the question it cannot answer about itself.

The distinction has a testable shape, because different audit responses to a
prior finding leave different marks on the outcomes of
Table~\ref{tab:fourstate}. Table~\ref{tab:multitype} works through three ways
an auditor might let last year's finding change this year's work, with the
underlying problems held identical and drawn independently each period, so
nothing underneath persists at all. Every cell is the change in the
prior-finding gap relative to a baseline in which prior findings change
nothing. The regimes share their first-period histories, so the history
groups and each program's previously flagged conditions and areas are the
same in every regime; they share second-period conditions and detection draws
program by program; and each regime changes detection only for the conditions
it targets. An outcome outside a regime's target set is therefore identical,
program by program, to its baseline value, which makes the zeros exact and
cancels the opportunity-set artifact of Section~\ref{sec:fac}. New findings are
separated by whether they fall in a requirement area that already had one,
because that is where the three responses part company.

\emph{Condition-specific follow-up.} If the response is to re-test the control
or re-examine the transactions behind the specific finding, repeat findings
rise and nothing else moves; the two new-finding outcomes are exactly zero.

\emph{Category-level follow-up.} If the response is to tighten work across the
whole requirement area that produced the finding, repeats rise and so do new
findings inside that area, while areas never flagged remain untouched.

\emph{Generalized scrutiny.} If the response is general wariness, so that a
program with a finding simply gets a harder audit everywhere, all three
outcomes move, including findings in requirement areas that never had one.
That last column is the one to watch, because a finding in a previously clean
area is exactly what an auditor would read as a genuinely new and independent
problem.

A new finding is therefore not automatically independent evidence of a new
underlying problem, because a prior finding may have changed how broadly the
auditor searched. The converse holds with equal force: broadly elevated
findings are also consistent with programs that are genuinely worse across the
board, and with stable differences between auditees that make some of them
durably more prone to deficiencies (Appendix~\ref{app:id}). The table narrows the space of explanations; it does not identify the
mechanism. In the Clearinghouse data the elevation runs across repeat-only,
mixed and any-new outcomes, which is the pattern the last two rows of
Table~\ref{tab:multitype} produce and the pattern a worse program produces,
and the record cannot say which.

The design question that follows is not how to squeeze more out of the
finding record but what additional information would separate the two factors
in equation~\eqref{eq:qrd}. Section~\ref{sec:scoring} shows what happens when
that question is skipped; Section~\ref{sec:operational} answers it.

\begin{table}[tbp]
\centering
\caption{\textbf{How you look decides what turns up.} Three ways of letting last
year's finding change this year's audit, with the underlying problems held
identical and drawn independently each period. Each cell is the change in the
gap between programs with and without a prior finding, measured against a
baseline in which prior findings change nothing. Following up the same finding
moves repeats and \emph{only} repeats. Looking harder across a requirement area
also surfaces new findings inside that area. Looking harder at the program
overall also surfaces findings in areas that never had one --- which is what an
auditor would otherwise read as a genuinely new problem. The zeros are exact:
they come from using the same random draws across regimes.}
\label{tab:multitype}
\small
\begin{tabular}{lccc}
\toprule
& \multicolumn{3}{c}{Change in the prior-finding gap} \\
\cmidrule(lr){2-4}
If last year's finding makes the auditor\dots
& \shortstack{A finding\\repeats}
& \shortstack{Something new,\\in an area that\\already had one}
& \shortstack{Something new,\\in an area that\\did not} \\
\midrule
Follow up the same finding & $+0.0394$ & $+0.0000$ & $+0.0000$ \\
Look harder at the same requirement area & $+0.0279$ & $+0.0454$ & $+0.0000$ \\
Look harder at the program overall & $+0.0223$ & $+0.0366$ & $+0.1388$ \\
\bottomrule
\end{tabular}
\end{table}

